% Part: sets-functions-relations
% Chapter: functions
% Section: isomorphisms

\documentclass[../../../include/open-logic-section]{subfiles}


\begin{document}

\olfileid[hi]{sfr}{fun}{iso}
\olsection{समरूपता}

\begin{explain}
\emph{समरूपता} ऐसा एकैकी आच्छादक प्रतिचित्रण है जो उन समुच्चयों की संरचना को
सुरक्षित रखता है जिन्हें वह परस्पर संबद्ध करता है; यहाँ संरचना का अर्थ उन
संबंधों से है जो समुच्चयों के !!{element}ों के बीच विद्यमान हों। निम्न दो
समुच्चयों $X=\{1,2,3\}$ और $Y=\{4,5,6\}$ पर विचार कीजिए। दोनों समुच्चयों को
परवर्ती, लघुतर और बृहत्तर संबंध संरचना प्रदान करते हैं। इन समुच्चयों के बीच
समरूपता ऐसा !!a{bijection} है जो इन संरचनाओं को सुरक्षित रखे। अतः !!a{bijective}
फलन $f \colon X
\to Y$ तब समरूपता है जब $i<j$ तभी और केवल तभी जब
$f(i)<f(j)$, $i>j$ तभी और केवल तभी जब $f(i)>f(j)$, और $j$,~$i$ का परवर्ती
तभी और केवल तभी हो जब $f(j)$,~$f(i)$ का परवर्ती हो।
\end{explain}

\begin{defn}[समरूपता]
मान लीजिए कि $U$ युग्म $\langle X, R\rangle$ और $V$ युग्म $\langle
Y, S\rangle$ हैं, जहाँ $X$ और $Y$ समुच्चय हैं तथा $R$ और $S$ क्रमशः $X$ और
$Y$ पर संबंध हैं। कोई !!{bijection} $f$,~$X$ से $Y$ तक,~$U$ से $V$ तक
\emph{समरूपता} तभी और केवल तभी है जब वह संबंधात्मक संरचना को सुरक्षित रखे;
अर्थात्, किन्हीं $x_{1}$ और $x_{2}$, जो $X$ में हों, के लिए,
$\tuple{x_1,x_2} \in R$ तभी और केवल तभी जब
$\tuple{f(x_1),f(x_2)} \in S$।
\end{defn}

\begin{ex}
निम्न दो समुच्चयों $X=\{1,2,3\}$ और $Y=\{4,5,6\}$ तथा लघुतर और बृहत्तर
संबंधों पर विचार कीजिए। फलन $f\colon X \to
Y$, जहाँ $f(x) = 7-x$ है, $\tuple{X,<}$ और $\tuple{Y,>}$ के बीच समरूपता
है।
\end{ex}

\end{document}
